% !TEX program = lualatex
% Generated by new_lecture_note.sh
\documentclass[11pt]{article}

\usepackage{fontspec}
\usepackage{microtype}
\usepackage{mathtools,amssymb,amsfonts,amsthm,bm}
\usepackage{enumitem}
\usepackage{booktabs,array,xcolor}
\usepackage{geometry,fancyhdr,setspace}
\usepackage[hidelinks]{hyperref}
\usepackage[nameinlink,noabbrev]{cleveref}

\geometry{margin=1in,headheight=15pt}
\setstretch{1.12}
\setlength{\parindent}{0pt}
\setlength{\parskip}{0.55em plus 0.12em minus 0.08em}
\setlength{\abovedisplayskip}{0.8em plus 0.2em minus 0.1em}
\setlength{\belowdisplayskip}{0.8em plus 0.2em minus 0.1em}
\allowdisplaybreaks[3]
\setlist{leftmargin=*,topsep=0.35em,itemsep=0.15em,parsep=0pt}

% ---------- Metadata ----------
\newcommand{\studentname}{Keshav Ramamurthy}
\newcommand{\coursename}{Math 104}
\newcommand{\lecturenumber}{01}
\newcommand{\lecturetitle}{Lecture 01}
\newcommand{\lecturedate}{\today}

% ---------- Reusable notation ----------
\newcommand{\N}{\mathbb{N}}
\newcommand{\Z}{\mathbb{Z}}
\newcommand{\Q}{\mathbb{Q}}
\newcommand{\R}{\mathbb{R}}
\newcommand{\C}{\mathbb{C}}
\newcommand{\F}{\mathbb{F}}
\newcommand{\K}{\mathbb{K}}
\newcommand{\E}{\mathbb{E}}
\newcommand{\Prob}{\mathbb{P}}
\newcommand{\1}{\mathbf{1}}
\newcommand{\eps}{\varepsilon}
\newcommand{\dd}{\,\mathrm{d}}
\newcommand{\abs}[1]{\left\lvert#1\right\rvert}
\newcommand{\norm}[1]{\left\lVert#1\right\rVert}
\newcommand{\ip}[2]{\left\langle#1,#2\right\rangle}
\newcommand{\set}[1]{\left\{#1\right\}}
\newcommand{\given}{\,\middle\vert\,}
\newcommand{\indep}{\perp\!\!\!\perp}
\newcommand{\lto}{\longrightarrow}
\newcommand{\deriv}[2]{\frac{\mathrm{d}#1}{\mathrm{d}#2}}
\newcommand{\pderiv}[2]{\frac{\partial#1}{\partial#2}}
\DeclareMathOperator{\Aut}{Aut}
\DeclareMathOperator{\Hom}{Hom}
\DeclareMathOperator{\Ker}{ker}
\DeclareMathOperator{\im}{im}
\DeclareMathOperator{\rank}{rank}
\DeclareMathOperator{\nullity}{nullity}
\DeclareMathOperator{\Span}{span}
\DeclareMathOperator{\Var}{Var}
\DeclareMathOperator{\Cov}{Cov}

% ---------- Note environments ----------
\newtheoremstyle{notestyle}{0.7em}{0.7em}{\itshape}{}{}{\bfseries}{.5em}{}
\theoremstyle{notestyle}
\newtheorem{theorem}{Theorem}[section]
\newtheorem{lemma}[theorem]{Lemma}
\newtheorem{proposition}[theorem]{Proposition}
\newtheorem{corollary}[theorem]{Corollary}
\theoremstyle{definition}
\newtheorem{definition}[theorem]{Definition}
\newtheorem{example}[theorem]{Example}
\newtheorem{remark}[theorem]{Remark}
\newenvironment{claim}{\par\noindent\textbf{Claim.}\ }{\par}
\newenvironment{idea}{\par\noindent\textit{Idea.}\ }{\par}
\newenvironment{proofsketch}{\par\noindent\textbf{Proof sketch.}\ }{\par}

% ---------- Header ----------
\pagestyle{fancy}
\fancyhf{}
\fancyhead[L]{\coursename}
\fancyhead[C]{\lecturetitle}
\fancyhead[R]{\studentname}
\fancyfoot[C]{\thepage}

\title{\vspace{-1.5em}\textbf{\coursename -- \lecturetitle}}
\author{\studentname}
\date{\lecturedate}

\begin{document}
\maketitle
\section{START}
\begin{definition}[Sequence]
\label{def:label}
A sequence is a map from ${n \in \mathbb{Z} | n \ge m}(m \in \mathbb{Z})$ to a set $X$
\end{definition}
\begin{definition}[Convergence]
\label{def:label}
A sequence $(s_n)_{1}^{\infty} \subset \mathbb{R} $ converges to $s \in \mathbb{R}$ if
$$\forall \epsilon > 0, \exists N > 0 \text{ such that }  n > N \implies |s_n - s| <
\epsilon$$
% \end{definition}
\begin{definition}[Divergence]
\label{def:label}
A sequence $(s_n)$ diverges if $(s_n)$ doesn't converge to any real number.
\end{definition}
\end{definition}
\begin{definition}[Limit Notation]
\label{def:label}
We write $$\lim_{n\rightarrow \infty}{(s_n)} = s, \lim s_n = s, s_n^{n \rightarrow \infty} \rightarrow s$$
We write $$lim_{n\rightarrow \infty}{(s_n)} = \infty$$
if $$\forall m \in \mathbb{R} \exists N > 0 \text { such that } n > N \implies s_n > M$$
\end{definition}
\begin{example}[Convergence Example]
\label{ex:label}
Take $s_n = \frac{3n+1}{7n+4}, n \in \mathbb{N}.$
We have 
\begin{align*}
    \lim_{n\rightarrow\infty}{(s_n)}
    &= \lim_{n\rightarrow\infty}{\frac{3n+1}{7n+4}}\\
    &= \lim_{n\rightarrow\infty}{\frac{3 + \frac{1}{n}}{7 + \frac{4}{n}}}\\
    &=  \frac{3}{7}\\
\end{align*}
How do we verify this?
\end{example}
\begin{proof}
    Epsilon Delta: Let $\epsilon > 0.$ Take $N = \frac{19}{49\epsilon}+\frac{4}{7}.$
    if $n > N \rightarrow |s_n - \frac{3}{7}| < \epsilon$
    \par The real step is finding $N$ with equation $$\left|
    \frac{3n+1}{7n+4}-\frac{3}{7}\right| < \epsilon$$
\end{proof}
\begin{theorem}[Uniqueness of the limit]
\label{thm:label}
"Any sequence has a unique limit"
\par Let $(s_n)\subset R.$ If $\lim_{n\rightarrow\infty}=s\in\mathbb{R}\text{ and }
\lim_{n\rightarrow \infty}{s_n}=t\in \mathbb{R} \text{ then } s = t$
\end{theorem}
\begin{proof}
    Establish
    $$\lim s_n = s\rightarrow \exists N_1 > 0 \text{ such that } n > N_1 \implies |s_n -
    s| < \frac{\epsilon}{2}$$
    \par Establish again $$\lim s_n = s\rightarrow \exists N_2 > 0 \text{ such that } n
    > N_2 \implies |s_n - s| < \frac{\epsilon}{2}$$ 
    Then $$\left|s-t\right| \le |s - s_n| + |s_n - t| <
    \frac{\epsilon}{2}+\frac{\epsilon}{2}  =\epsilon$$
    Take $N = \max {N_1, N_2}$ Note that $N > N_1$ and $N > N_2$
\end{proof}
\section{Main idea}
\textit{Write the one-sentence point of today's lecture here.}

\section{Definitions and notation}
\begin{definition}[First definition]
Write the definition and one sentence explaining why it matters.
\end{definition}

\section{Claims, theorems, and proof sketches}
\begin{claim}
State the claim in one precise sentence.
\end{claim}

\begin{proofsketch}
Record the key idea, the first useful equation, and the step that still needs checking.
\end{proofsketch}

\section{Examples and calculations}
\begin{example}[Lecture example]
Write the setup, the calculation, and the conclusion.
\end{example}

\section{Questions and follow-up}
\begin{itemize}
  \item What is still unclear?
  \item Which exercise or reading should I revisit?
  \item TODO: 
\end{itemize}

\end{document}
